% =====================================================================
% Abstract (en inglés, según la normativa del TFM)
% =====================================================================
\begin{otherlanguage}{english}
\chapter*{Abstract}
\addcontentsline{toc}{chapter}{Abstract}

Positron Emission Tomography (PET) exposes patients to ionising
radiation, which motivates acquiring scans at a reduced dose; doing so,
however, degrades image quality and turns the recovery of a full-dose
image from its low-dose counterpart into an ill-posed inverse problem.
This thesis studies denoising diffusion probabilistic models as a
generative tool for low-dose PET reconstruction and contributes a
comparative analysis of two distinct diffusion-based methods. The first
is a \emph{supervised conditional} pipeline, in which the low-dose scan
is fed as a conditioning channel to a UNet trained to reverse the
diffusion process. The second is an \emph{unconditional prior with
measurement guidance}, which learns a full-dose prior and injects the
observation only at inference time through Diffusion Posterior Sampling
(DPS). Both methods are built and evaluated on a dataset of 371 paired
NIfTI volumes (full dose and dose reduced to 1/20), operating slice by
slice over axial cross-sections. To delimit the contribution of the
generative paradigm, the two diffusion methods are benchmarked under
identical data, normalisation and training budget against two classical
discriminative baselines: a regression UNet (an ablation of the
diffusion process) and a RED-CNN (the reference convolutional
architecture for low-dose medical-image denoising). All models are
compared on the same held-out test set using image-quality metrics
(PSNR, SSIM, NRMSE) and intensity-preservation metrics. The results show
that the diffusion methods reconstruct anatomically faithful full-dose
images and substantially improve over the raw low-dose observation,
while the classical baselines remain slightly superior in the aggregate
metrics, a finding discussed together with the associated computational
cost.

\vspace{1em}
\noindent\textbf{Keywords:} diffusion models, PET reconstruction,
low-dose PET, inverse problems, generative AI, DDPM, DDIM, DPS,
medical imaging.

\end{otherlanguage}
